
\documentclass{article}
\usepackage{colortbl}
\usepackage{makecell}
\usepackage{multirow}
\usepackage{supertabular}

\begin{document}

\newcounter{utterance}

\centering \large Interaction Transcript for game `mmlu\_pro', experiment `mmlu\_pro\_chemistry', episode 1090 with Qwen3.6{-}27B.
\vspace{24pt}

{ \footnotesize  \setcounter{utterance}{1}
\setlength{\tabcolsep}{0pt}
\begin{supertabular}{c@{$\;$}|p{.15\linewidth}@{}p{.15\linewidth}p{.15\linewidth}p{.15\linewidth}p{.15\linewidth}p{.15\linewidth}}
    \# & \multicolumn{2}{c}{Player} && \multicolumn{2}{c}{Game Master} \\
    \hline

    \theutterance \stepcounter{utterance}  
    & & & \multicolumn{4}{p{0.6\linewidth}}{
        \cellcolor[rgb]{0.9,0.9,0.9}{
            \makecell[{{p{\linewidth}}}]{
                \texttt{\tiny{[P1$\langle$GM]}}
                \texttt{The following are multiple choice questions (with answers) about chemistry. Answer only with "The answer is (X)" where X is the correct letter choice.} \\
\\ 
\texttt{Question:} \\
\texttt{Predict the number of lines in the EPR spectrum of a solution of 13C{-}labelled methyl radical (13CH3•), assuming the lines do not overlap.} \\
\texttt{Options:} \\
\texttt{A. 10} \\
\texttt{B. 8} \\
\texttt{C. 4} \\
\texttt{D. 20} \\
\texttt{E. 12} \\
\texttt{F. 3} \\
\texttt{G. 16} \\
\texttt{H. 5} \\
\texttt{I. 24} \\
\texttt{J. 6} \\
\texttt{The answer is (B)} \\
\\ 
\texttt{Question:} \\
\texttt{Which of the following lists the hydrides of group{-}14 elements in order of thermal stability, from lowest to highest?} \\
\texttt{Options:} \\
\texttt{A. GeH4 < SnH4 < PbH4 < SiH4 < CH4} \\
\texttt{B. SiH4 < GeH4 < SnH4 < PbH4 < CH4} \\
\texttt{C. PbH4 < CH4 < SnH4 < GeH4 < SiH4} \\
\texttt{D. PbH4 < SnH4 < CH4 < GeH4 < SiH4} \\
\texttt{E. SnH4 < GeH4 < SiH4 < PbH4 < CH4} \\
\texttt{F. CH4 < GeH4 < SnH4 < PbH4 < SiH4} \\
\texttt{G. SiH4 < SnH4 < PbH4 < GeH4 < CH4} \\
\texttt{H. CH4 < SiH4 < GeH4 < SnH4 < PbH4} \\
\texttt{I. CH4 < PbH4 < GeH4 < SnH4 < SiH4} \\
\texttt{J. PbH4 < SnH4 < GeH4 < SiH4 < CH4} \\
\texttt{The answer is (J)} \\
\\ 
\texttt{Question:} \\
\texttt{Which of the following is considered an acid anhydride?} \\
\texttt{Options:} \\
\texttt{A. H2SO3} \\
\texttt{B. NaCl} \\
\texttt{C. CH4} \\
\texttt{D. HNO3} \\
\texttt{E. SO2} \\
\texttt{F. Al(NO3)3} \\
\texttt{G. CaCO3} \\
\texttt{H. H2O} \\
\texttt{I. HCl} \\
\texttt{J. N/A} \\
\texttt{The answer is (E)} \\
\\ 
\texttt{Question:} \\
\texttt{A new compound is synthesized and found to be a monoprotic acid with a molar mass of 248 g/mol. When 0.0050 mol of this acid are dissolved in 0.500 L of water, the pH is measured as 3.89. What is the pKa of this acid?} \\
\texttt{Options:} \\
\texttt{A. 5.78} \\
\texttt{B. 4.78} \\
\texttt{C. 4.56} \\
\texttt{D. 6.89} \\
\texttt{E. 7.78} \\
\texttt{F. 3.89} \\
\texttt{G. 1.23} \\
\texttt{H. 2.89} \\
\texttt{I. 2.33} \\
\texttt{J. 5.33} \\
\texttt{The answer is (A)} \\
\\ 
\texttt{Question:} \\
\texttt{A solution contains 2.00 mole of acetic acid, CH3COOH, and 1.00 mole of calcium acetate, Ca(CH3COO)2. The solution is able to resist the addition of a small amount of strong acid or strong base with only minor changes in the pH of the solution. Larger quantities of strong acid or strong base can cause a significant change in pH. How many moles of nitric acid, HNO3, may be added before the pH begins to change significantly?} \\
\texttt{Options:} \\
\texttt{A. 0.250 mole} \\
\texttt{B. 0.500 mole} \\
\texttt{C. 3.00 mole} \\
\texttt{D. 1.00 mole} \\
\texttt{E. 3.50 mole} \\
\texttt{F. 1.50 mole} \\
\texttt{G. 2.50 mole} \\
\texttt{H. 4.00 mole} \\
\texttt{I. 0.750 mole} \\
\texttt{J. 2.00 mole} \\
\texttt{The answer is (J)} \\
\\ 
\texttt{Question:} \\
\texttt{What volume of hydrogen at STP is produced as sulfuric acid acts on 120 g. of metallic calcium. Equation for the re{-}action is Ca + H\_2 SO\_4 \textbackslash{}rightarrow CaSO\_4 + H\_2} \\
\texttt{Options:} \\
\texttt{A. 33.6 liters} \\
\texttt{B. 56.0 liters} \\
\texttt{C. 3.00 liters} \\
\texttt{D. 84.0 liters} \\
\texttt{E. 67.2 liters} \\
\texttt{F. 44.8 liters} \\
\texttt{G. 134.4 liters} \\
\texttt{H. 120 liters} \\
\texttt{I. 11.2 liters} \\
\texttt{J. 22.4 liters} \\
            }
        }
    }
     \\ \\

    \theutterance \stepcounter{utterance}  
    & \multicolumn{4}{p{0.6\linewidth}}{
        \cellcolor[rgb]{0.561,0.737,0.561}{
            \makecell[{{p{\linewidth}}}]{
                \texttt{\tiny{[P1$\rangle$GM]}}
                \texttt{The answer is (F)} \\
            }
        }
    }
    & & \\ \\

    \theutterance \stepcounter{utterance}  
    & & & \multicolumn{2}{p{0.3\linewidth}}{
        \cellcolor[rgb]{0.9,0.9,0.9}{
            \makecell[{{p{\linewidth}}}]{
                \texttt{\tiny{[GM$|$GM]}}
                \texttt{F} \\
            }
        }
    }
    & & \\ \\

    \theutterance \stepcounter{utterance}  
    & & & \multicolumn{2}{p{0.3\linewidth}}{
        \cellcolor[rgb]{0.9,0.9,0.9}{
            \makecell[{{p{\linewidth}}}]{
                \texttt{\tiny{[GM$|$GM]}}
                \texttt{E} \\
            }
        }
    }
    & & \\ \\

    \theutterance \stepcounter{utterance}  
    & & & \multicolumn{2}{p{0.3\linewidth}}{
        \cellcolor[rgb]{0.9,0.9,0.9}{
            \makecell[{{p{\linewidth}}}]{
                \texttt{\tiny{[GM$|$GM]}}
                \texttt{game\_result = LOSE} \\
            }
        }
    }
    & & \\ \\

\end{supertabular}
}

\end{document}
